\section{灵敏度分析与模型检验}

\subsection{灵敏度分析方法}

为评估模型关键参数对优化结果的影响程度，选取购电电价、光伏装机容量、风电装机容量和售电电价四个参数进行灵敏度分析。以问题三典型场景下54吨/日产量的最优方案为基准（基准吨氨成本4530.88元/吨），分别对各参数进行$\pm$20\%--30\%的扰动，观察吨氨成本的变化幅度。

\subsection{灵敏度分析结果}

\begin{table}[H]
\centering
\caption{关键参数灵敏度分析结果}\label{tab:sensitivity}
\begin{tabular}{lcccc}
\toprule
参数 & 扰动范围 & 成本降低(元/吨) & 成本升高(元/吨) & 敏感度排名 \\
\midrule
购电电价 & $\pm$30\% & $-$501 & +501 & 1（最敏感） \\
光伏装机容量 & $\pm$20\% & $-$382 & +493 & 2 \\
风电装机容量 & $\pm$20\% & $-$259 & +266 & 3 \\
售电电价 & $\pm$30\% & $-$85 & +85 & 4（最不敏感） \\
\bottomrule
\end{tabular}
\end{table}

灵敏度分析结果（表\ref{tab:sensitivity}）揭示了以下规律：

购电电价是影响吨氨成本最敏感的参数，$\pm$30\%的扰动导致成本变化$\pm$501元/吨（相对变化$\pm$11.1\%）。这是因为园区在风光不足时段仍需大量购电，购电成本在总成本中占比最高。该结果对园区选址和电价谈判具有重要指导意义——优先选择低电价地区建设园区可显著降低生产成本。

光伏装机容量的敏感度排名第二，且呈现不对称特征：容量减少20\%导致成本升高493元/吨，而容量增加20\%仅降低成本382元/吨。这种不对称性源于光伏出力集中在白天，增加光伏容量在白天已经过剩的时段效果有限（多余电力只能低价上网），但减少光伏容量会迫使更多时段购电。

风电装机容量的敏感度排名第三，变化幅度相对对称（$\pm$260元/吨左右）。风电出力在全天分布较为均匀，增减容量对购电需求的影响较为线性。

售电电价是最不敏感的参数，$\pm$30\%扰动仅导致$\pm$85元/吨的成本变化。这是因为在优化调度下，园区尽量减少余电上网（以满足绿电指标），售电量本身较小，售电电价变化对总成本的影响有限。

\subsection{模型检验}

\subsubsection{能量守恒验证}

对所有计算场景进行能量守恒检验：$E_{RE} + E_{buy} = E_{total} + E_{sell}$。所有场景的能量平衡误差均小于0.001 MWh（相对误差$<$0.0002\%），验证了功率平衡模型的正确性。

\subsubsection{约束满足验证}

检验所有优化结果是否满足物理约束：（1）所有功率变量非负；（2）$\alpha(t) \in [0.1, 1.0]$（问题三、四）；（3）购售电互斥（同一时段$P_{buy}(t) \cdot P_{sell}(t) = 0$）；（4）储能SOC在$[0, C_{sto}]$范围内；（5）日产量不超过72吨/日。所有约束均严格满足。

\subsubsection{结果单调性验证}

验证结果的物理合理性：（1）连续调节成本$\leq$离散调节成本（更灵活的调度空间不会使结果变差）——验证通过，4209.48 $<$ 4493.82；（2）有储能产量$\geq$无储能产量——验证通过，10619 $>$ 9782；（3）联网成本$\leq$离网成本——验证通过，4209.48 $<$ 4785.85。所有单调性关系均符合物理直觉。
